\documentclass{article}
\usepackage[T1]{fontenc}
\usepackage{lmodern}
\usepackage{graphicx}
\usepackage{amsmath,amssymb}
\usepackage[a4paper,margin=2cm]{geometry}
\usepackage[absolute,overlay]{textpos}
\setlength{\TPHorizModule}{1mm}
\setlength{\TPVertModule}{1mm}
\usepackage[indonesian]{babel}
\usepackage{hyperref}
\setlength{\emergencystretch}{2em}
% unit: Halaman 69

\begin{document}

% === Halaman 69 (llm) ===

\begin{enumerate}
  \item[3.] Ketiga register didetak selama 100 putaran dalam fase $stop/go$ dengan menggunakan kaidah mayoritas, namun bit-bit luarannya dibuang (tidak dipakai).
  \item[4.] Selanjutnya, ketiga register didetak selama 228 putaran dalam fase $stop/go$ menggunakan kaidah mayoritas untuk menghasilkan bit-bit kunci-alir sepanjang 228 bit.
\end{enumerate}

Pada setiap putaran dihasilkan 1 bit yang merupakan hasil peng-XOR-an bit-bit MSB yang terlempar. Total bit luaran adalah 228 bit. Kunci-alir 228-bit inilah yang kemudian digunakan untuk mengenkripsi $frame$ pesan sepanjang 228 bit.

\end{document}
